\section*{ВВЕДЕНИЕ}
\addcontentsline{toc}{section}{ВВЕДЕНИЕ}

Интерактивные приложения занимают важное место среди современных программных систем. К ним относятся учебные и инженерные симуляторы, демонстрационные трехмерные сцены, визуализаторы, игровые прототипы, виртуальные лаборатории и другие программы, в которых результат работы зависит от постоянного взаимодействия пользователя с объектами сцены. Такие приложения должны не только выводить изображение на экран, но и регулярно обновлять внутреннее состояние, обрабатывать ввод с клавиатуры и мыши, управлять временем кадра, изменять положение объектов, выполнять пользовательские сценарии поведения и формировать новое изображение с малой задержкой.

Разработка подобных программ связана с повторяющимся набором технических задач. Разработчику необходимо создать окно приложения, организовать основной цикл, хранить объекты сцены, задавать их свойства, подключать графические ресурсы, описывать поведение объектов, обрабатывать пользовательский ввод и, при необходимости, рассчитывать простые физические взаимодействия. Если все эти механизмы реализуются заново для каждого отдельного приложения, значительная часть времени тратится не на прикладную логику, а на повторное создание базовой программной инфраструктуры.

Для решения этой проблемы используются программные основы и движки, предоставляющие готовые механизмы работы со сценами, объектами, ресурсами и графикой. Однако существующие решения не всегда удобны для небольших учебных и исследовательских проектов. Полноценные игровые движки и крупные фреймворки обладают большим количеством возможностей, но могут быть избыточными, сложными в изучении и не всегда позволяют подробно рассмотреть внутреннее устройство системы. Низкоуровневые библиотеки, напротив, дают доступ к графическому API и средствам ввода, но оставляют разработчику необходимость самостоятельно проектировать архитектуру приложения. Например, LWJGL предоставляет Java-доступ к низкоуровневым библиотекам, в том числе OpenGL и GLFW, но не задает готовую объектную модель приложения \cite{lwjgl}. OpenGL, в свою очередь, решает задачу вывода графики, но не определяет способ организации сцен, компонентов и систем обработки \cite{khronos-opengl}.

В связи с этим актуальной является разработка компонентно-ориентированной программной системы, которая занимает промежуточное положение между низкоуровневой графической библиотекой и крупным готовым движком. Такая система должна предоставлять Java-разработчику основные средства построения интерактивного приложения, но при этом оставаться достаточно простой и открытой для изучения, модификации и расширения.

Основой разрабатываемой системы является компонентный подход к организации объектов сцены. В рамках этого подхода объект не описывается глубокой иерархией наследования. Вместо этого он представляется сущностью, к которой добавляются независимые компоненты: трансформация, графическое представление, камера, источник света, физическое тело, коллайдер или пользовательский скрипт. Обработка таких объектов выполняется системами, каждая из которых отвечает за отдельную часть поведения приложения. Подобная организация позволяет переиспользовать компоненты, добавлять новую функциональность без изменения базового класса объекта и разделять разные аспекты работы сцены.

Разрабатываемая программная система ориентирована на создание настольных интерактивных приложений на языке Java. Ее пользователем является не конечный пользователь готового приложения, а разработчик, который создает классы сцен, добавляет сущности и компоненты, подключает системы обработки и описывает поведение объектов средствами программного кода. В составе системы рассматриваются ядро компонентной модели, управление приложением и сценами, графическая подсистема, подсистема ресурсов, обработка пользовательского ввода, скрипты, отладочная визуализация и базовая физическая подсистема.

\emph{Цель настоящей работы} -- разработка компонентно-ориентированной программной системы для построения интерактивных приложений на Java.

Для достижения поставленной цели необходимо решить следующие задачи:
\begin{itemize}
\item провести анализ предметной области и существующих программных средств построения интерактивных приложений;
\item сформулировать требования к компонентно-ориентированной программной системе и определить порядок ее использования разработчиком;
\item спроектировать архитектуру системы, включающую ядро компонентной модели, управление сценами, графическую подсистему, обработку ввода, скрипты и базовую физику;
\item реализовать основные классы, компоненты, системы обработки и демонстрационные сцены;
\item выполнить модульное, регрессионное и системное тестирование разработанной программной системы.
\end{itemize}

\emph{Структура работы.} Выпускная квалификационная работа состоит из введения, четырех разделов основной части, заключения, списка использованных источников и приложений.

\emph{Во введении} обоснована актуальность выбранной темы, сформулированы цель и задачи разработки, а также приведена структура выпускной квалификационной работы.

\emph{В первом разделе} выполняется анализ предметной области. Рассматриваются интерактивные приложения как класс программных систем, особенности их внутренней организации, компонентно-ориентированный подход, ключевые сущности предметной области, особенности разработки на языке Java и существующие программные решения.

\emph{Во втором разделе} приводится техническое задание. В нем указано основание для разработки, определены цель и назначение программной системы, сформулированы требования к архитектурной модели, управлению приложением и сценами, графической подсистеме, пользовательскому вводу, скриптам, базовой физике, демонстрационным сценам, программному интерфейсу, нефункциональным характеристикам и оформлению документации.

\emph{В третьем разделе} представлен технический проект программной системы. В разделе обоснован выбор технологий, описаны Java, LWJGL, OpenGL, GLSL, GLFW и STB, рассмотрено внутреннее устройство системы, структура пакетов, жизненный цикл приложения и сцены, а также проектирование основных подсистем и демонстрационных сцен.

\emph{В четвертом разделе} приведен рабочий проект. В нем представлена спецификация реализованных классов и компонентов программной системы, описаны математические, графические, ресурсные и физические классы, системы обработки и демонстрационные сцены. Также в разделе приведены результаты модульного, регрессионного и системного тестирования.

В заключении излагаются основные результаты, полученные при выполнении выпускной квалификационной работы.

В приложении А представлен графический материал, подготовленный для защиты выпускной квалификационной работы. В приложении Б приведены фрагменты исходного кода разработанной программной системы.
